\section*{ЗАКЛЮЧЕНИЕ}
\addcontentsline{toc}{section}{ЗАКЛЮЧЕНИЕ}

В выпускной квалификационной работе рассмотрена разработка клиентской подсистемы программно-информационной системы <<Fitness Life>> -- бизнес-проекта стартап-команды, направленного на индивидуальное планирование тренировок, учёт питания и сопутствующую мотивацию пользователя в едином веб-приложении. Выполненный инженерный вклад автора относится к браузерной части продукта: проектированию одностраничного интерфейса, модульной организации клиентского кода и обеспечению взаимодействия с REST API серверной подсистемы, реализуемой напарником по проекту.

В ходе работы изучена предметная область индивидуальных занятий физической культурой, проанализированы аналогичные цифровые сервисы и сформулированы требования к клиентскому приложению. На основе анализа обоснована клиент-серверная архитектура: серверная часть отвечает за аутентификацию, хранение профилей и синхронизируемого пользовательского хранилища, а клиентская -- за представление данных, ввод пользователем и выполнение прикладных сценариев без полной перезагрузки страницы.

Спроектирована и реализована клиентская подсистема в каталоге \texttt{ДИПЛОМ 2} репозитория проекта. Приложение построено как одностраничное веб-приложение (HTML, CSS, JavaScript): единая HTML-страница задаёт каркас экранов авторизации и основной оболочки, модуль \texttt{auth.js} обеспечивает регистрацию, вход, работу с JWT и синхронизацию \texttt{localStorage} с сервером, модуль \texttt{app.js} -- логику тренировочных планов, журнала подходов и визуализации нагрузки с использованием Chart.js. Предметные разделы (профиль, питание, поддержка, социальная лента, аватар) вынесены в отдельные файлы \texttt{*.module.js}, что упрощает сопровождение и развитие функциональности.

Основные результаты работы:
\begin{enumerate}
\item Проведён анализ предметной области и сопоставлены подходы к построению клиентских веб-приложений в сфере фитнеса; определены функциональные требования к клиентской части системы.
\item Разработано техническое задание и технический проект клиентской подсистемы: описаны архитектура одностраничного приложения, состав модулей, работа с DOM, Web Storage API и событийной моделью браузера.
\item Реализованы экраны аутентификации и профиля, клиентские модули тренировок, питания, психологической поддержки, цифрового аватара и социальной ленты.
\item Организовано локальное хранение рабочего состояния пользователя и синхронизация данных с REST API серверной части при наличии сетевого соединения.
\item Выполнено системное тестирование ключевых пользовательских сценариев в браузере: регистрация и вход, работа с планами и журналом тренировок, отображение графиков прогресса, ведение дневника питания, публикации в ленте, чекины в разделе поддержки, отображение аватара и редактирование профиля.
\end{enumerate}

Поставленные во введении и в задании на выпускную квалификационную работу задачи решены. Требования технического задания к клиентской подсистеме выполнены в объёме, необходимом для демонстрации MVP продукта <<Fitness Life>> и согласованной работы с серверной частью. Готовый рабочий проект представлен исходным кодом клиентских модулей, спецификацией методов и результатами проверки интерфейса; фрагменты кода и графический материал приведены в приложениях.

Направления дальнейшего развития клиентской подсистемы включают: замену заглушек скриншотов в разделе тестирования на иллюстрации реального интерфейса; повышение доступности (ARIA, навигация с клавиатуры); оптимизацию загрузки (разбиение на бандлы, ленивая подгрузка модулей); расширение офлайн-режима за счёт Service Worker; углубление аналитики тренировок и интеграцию с носимыми устройствами при развитии серверного API; внедрение элементов freemium-модели монетизации, обозначенной в бизнес-обосновании проекта.

Таким образом, в работе создана работоспособная клиентская часть веб-приложения <<Fitness Life>>, обеспечивающая пользователю единый интерфейс для планирования нагрузки, учёта питания и социально-мотивационных функций с сохранением персональных данных и возможностью их синхронизации между сеансами работы в браузере.
